\begin{trapbox}
  \textbf{\textcolor{CTrap}{易错点一｜加在括号外，还是括号内？}}
  \textbf{错法：}把$y=x^2$上移$1$，写成$y=(x+1)^2$。

  \textbf{正确：}上移$1$写$\boxed{y=x^2+1}$；$y=(x+1)^2$其实是\textbf{左移$1$}。前者改高度，后者改位置。
  \begin{center}
  \begin{tikzpicture}[x=0.96cm,y=0.69cm,>=stealth]
    \draw[->,gray] (-2.35,0)--(2.15,0) node[right] {$x$};
    \draw[->,gray] (0,-0.45)--(0,3.50) node[above] {$y$};
    \draw[gray!80,thick,domain=-1.65:1.65,samples=55] plot (\x,{\x*\x});
    \draw[CTrap,very thick,domain=-1.45:1.45,samples=55] plot (\x,{\x*\x+1});
    \draw[blue!70!black,thick,dashed,domain=-2.60:0.60,samples=55] plot (\x,{(\x+1)*(\x+1)});
    \fill (0,0) circle (1.7pt);
    \fill[CTrap] (0,1) circle (1.7pt);
    \fill[blue!70!black] (-1,0) circle (1.7pt);
    \draw[->,CTrap,thick] (0.20,0.15)--(0.20,0.85);
    \draw[->,blue!70!black,thick] (-0.16,-0.24)--(-0.90,-0.24);
    \node[right,CTrap] at (2.0,3.4) {$x^2+1$};
    \node[left,blue!70!black] at (-2.9,1.6) {$(x+1)^2$};
    \node[right,gray!80!black] at (2.0,1.65) {$x^2$};
  \end{tikzpicture}

  \end{center}
  \textbf{一秒自查：}如果是\textbf{上下平移}，原来$x$能取哪些值，就仍能取哪些值；式子外面加减才对。

  \medskip
  \textbf{\textcolor{CTrap}{易错点二｜距离不能是负数；负号代表方向}}
  \textbf{易混说法：}“向下平移$-2$个单位”不规范，因为\textbf{距离不能写成负数}，同时又写“向下”和“$-2$”容易误解。

  \textbf{正确：}只用正距离表述：\textbf{向下移$2$个单位}，新式$y=f(x)-2$。如果题目明确使用\textbf{带符号的变化量}$t=-2$，它也表示\textbf{向下移$2$}，并不是向上移。
  \begin{center}
  \begin{tikzpicture}[x=0.95cm,y=0.62cm,>=stealth]
    \draw[->,gray] (0,-0.30)--(0,4.50) node[above] {$y$};
    \draw[gray!50,dashed] (0,3.40)--(2.70,3.40);
    \draw[gray!50,dashed] (0,1.40)--(2.70,1.40);
    \fill[gray!85] (1.40,3.40) circle (2pt) node[right] {原高度$3$};
    \fill[CTrap] (1.40,1.40) circle (2pt) node[right] {新高度$1$};
    \draw[->,CTrap,very thick] (1.05,3.25)--(1.05,1.55) node[midway,left] {$t=-2$};
    \node[left] at (0,3.40) {$3$};
    \node[left] at (0,1.40) {$1$};
  \end{tikzpicture}

  \end{center}
  \textbf{记住：}$b>0$是距离，$t\in\mathbb R$是带方向的变化量，两种记法不要混用。

  \medskip
  \textbf{\textcolor{CTrap}{易错点三｜上下平移不会改变定义域}}
  \textbf{错法：}把$y=\sqrt{x}$上移$1$，写成$y=\sqrt{x+1}+1$，误以为$x$的范围也要变化。

  \textbf{正确：}$\boxed{y=\sqrt{x}+1}$，仍然要有$x\ge0$；因为上下移动只改变每个点的高度，\textbf{不会产生原来没有的横坐标}。
  \begin{center}
  \begin{tikzpicture}[x=0.82cm,y=1.00cm,>=stealth]
    \draw[->,gray] (-0.70,0)--(4.60,0) node[right] {$x$};
    \draw[->,gray] (0,-0.30)--(0,3.35) node[above] {$y$};
    \draw[gray!80,thick,domain=0:4.0,samples=65] plot (\x,{sqrt(\x)});
    \draw[CTrap,very thick,domain=0:4.0,samples=65] plot (\x,{sqrt(\x)+1});
    \fill[gray!85] (0,0) circle (1.8pt);
    \fill[CTrap] (0,1) circle (1.8pt);
    \draw[->,CTrap,thick] (0.23,0.14)--(0.23,0.86) node[midway,right] {上移$1$};
    \node[right,gray!85!black] at (4.25,1.6) {$y=\sqrt{x}$};
    \node[right,CTrap] at (4.25,2.95) {$y=\sqrt{x}+1$};
    \node[below right] at (-.5,-.4) {$x=0$};
    \draw[<->,gray!70] (0,-.85)--(4.00,-.85);
    \node[below] at (2.10,-.85) {$x\ge0$};
  \end{tikzpicture}

  \end{center}
  \textbf{一秒自查：}上下移动前后，每个点的\textbf{横坐标原样保留}；原函数有什么取值限制，新函数就保留什么限制。
\end{trapbox}
